% Source: OpenAI2026AreaLaw, Lemma 7.2, 06-transport.tex, lines 157--172.
% Manuscript revision: adc7f1241b42e322a6451854ab7e4b4c146bf78a.
\section{Hermitian intertwining and matrix sandwiches}
\label{sec:matrix_sandwich_intertwine}

The following algebraic identities are auxiliary to the sandwich formulas
in the proof of Lemma~7.2 of~\cite{OpenAI2026AreaLaw}.
They hold for finite matrices over a semiring with involution; a
multiplicative identity is not required. The notation $L^*$ denotes the
conjugate transpose, and a matrix is Hermitian when $L^*=L$.

\begin{theorem}[The adjoint intertwining identity]
  \label{thm:matrix_adjoint_intertwine}
  \lean{Matrix.IsHermitian.conjTranspose_intertwine}
  \leanok
  Let $L_1,L_2$ be Hermitian square matrices, and let $J$ be a rectangular
  matrix satisfying $L_1J=JL_2$. Then
  \begin{align}
    L_2J^*&=J^*L_1.
  \end{align}
\end{theorem}

\begin{proof}
  \leanok
  Taking adjoints gives $J^*L_1^*=L_2^*J^*$.
  The two Hermitian identities give the conclusion.
\end{proof}

\begin{theorem}[Intertwining a sandwich]
  \label{thm:matrix_sandwich_intertwine}
  \lean{Matrix.IsHermitian.sandwich_intertwine}
  \leanok
  Under the hypotheses of Theorem~\ref{thm:matrix_adjoint_intertwine},
  every square matrix $H$ of the same size as $L_2$ satisfies
  \begin{align}
    L_1(JHJ^*)L_1&=J(L_2HL_2)J^*.
  \end{align}
  No Hermitian assumption on $H$ or isometry assumption on $J$ is required.
\end{theorem}

\begin{proof}
  \leanok
  \uses{thm:matrix_adjoint_intertwine}
  Associativity, the original intertwining identity, and
  Theorem~\ref{thm:matrix_adjoint_intertwine} give
  \begin{align}
    L_1(JHJ^*)L_1
      &=(L_1J)H(J^*L_1)
       =(JL_2)H(L_2J^*)
       =J(L_2HL_2)J^*.
  \end{align}
\end{proof}
